% ============================================
% Instructional Professor Cover Letter – University of Chicago
% ============================================

\documentclass[11pt,letterpaper]{letter}

\usepackage{geometry}
\usepackage{microtype}
\usepackage[hidelinks]{hyperref}

\geometry{left=25mm,right=25mm,top=25mm,bottom=25mm}

\signature{Decory Edwards}
\address{
Department of Economics \\
Johns Hopkins University \\
Baltimore, MD 21218 \\
\texttt{dedwar65@jh.edu}
}

\begin{document}

\begin{letter}{Search Committee \\
Kenneth C. Griffin Department of Economics \\
The University of Chicago \\
Chicago, IL}

\opening{Dear Members of the Search Committee,}

I am writing to apply for the Instructional Professor position in the Kenneth C. Griffin Department of Economics at the University of Chicago. I am a Ph.D.\ candidate in Economics at Johns Hopkins University, advised by Professors Christopher D.~Carroll, Nicholas W.~Papageorge, and Stelios Fourakis, and I expect to receive my degree in May 2026. My research fields are macroeconomics, wealth and income dynamics, and computational economics.

My research agenda focuses on understanding the determinants of wealth inequality and the mechanisms through which heterogeneity in returns to wealth shapes aggregate outcomes. In my job-market paper, I estimate the degree of heterogeneity in individual portfolio returns using micro-data from the Survey of Consumer Finances and simulate a structural model in which households face idiosyncratic return risk. I show that allowing for persistent heterogeneity in returns substantially improves the model’s ability to match the empirical distribution of wealth, offering a tractable way to connect micro-level return dispersion to macroeconomic inequality.

In a second project, I use the Health and Retirement Study to examine the relationship between interpersonal trust and portfolio performance. Building on the literature that finds a hump-shaped relationship between trust and income, I show that a similar relationship holds for individual returns to wealth measured in the HRS data, highlighting a behavioral channel through which social attitudes shape saving and investment outcomes. This work also provides new estimates of the persistent component of returns to net worth, which motivate the calibration of heterogeneity in my structural model.

The instructional mission of the University of Chicago’s Economics Department strongly aligns with my interests and experience. I am particularly drawn to a career-track teaching role that emphasizes excellence in undergraduate and master’s-level instruction, close student advising, and active involvement in curriculum development. The department’s commitment to rigorous economic reasoning, clarity of exposition, and analytical precision closely matches my own approach to teaching economics.

\textbf{Teaching.} I have experience teaching and supporting courses in \textit{Principles of Macroeconomics}, \textit{Intermediate Macroeconomic Theory}, and quantitative economics. My teaching emphasizes clear economic intuition, disciplined use of models, and the integration of data to reinforce theoretical concepts. I place particular emphasis on helping students develop economic reasoning skills that translate across fields, whether they pursue further study, policy work, or private-sector careers. At Chicago, I would be well prepared to teach a range of undergraduate and master’s courses, to mentor students with diverse academic backgrounds, and to contribute to advising, student activities, and curriculum coordination under the direction of the department’s instructional leadership.

The University of Chicago’s longstanding emphasis on serious engagement with economics, combined with the Griffin Department’s commitment to high-quality teaching at scale, makes this position especially appealing to me. I would welcome the opportunity to contribute to a department that values instructional excellence as a central academic mission and to support students as they develop strong analytical foundations in economics.

Thank you very much for your time and consideration. I would be grateful for the opportunity to discuss my application further.

\closing{Sincerely,}

\end{letter}
\end{document}
